\section*{Введение}
\addcontentsline{toc}{section}{Введение}

Современные вычислительные системы всё чаще проектируются как распределённые кластеры, в которых данные должны быть доступны одновременно с множества узлов, сохраняться при отказе отдельных компонентов и легко масштабироваться. Традиционные файловые системы, рассчитанные на один сервер, в таких условиях оказываются неэффективными и неотказоустойчивыми. На смену им приходят распределённые файловые системы, одной из которых является GlusterFS.

GlusterFS~--- это масштабируемая сетевая файловая система с открытым исходным кодом, предназначенная для объединения дисковых ресурсов множества серверов в единое логическое хранилище. Она применяется в задачах высокопроизводительных вычислений, облачных инфраструктурах и системах хранения больших данных благодаря поддержке репликации, распределения данных, квотирования и снимков состояния.

Работа выполнялась в VirtualBox на macOS (MacBook Air, процессор Apple Silicon). На всех узлах установлена ОС Debian~12 (архитектура arm64) и GlusterFS версии~10.3.
